\section{The first mixed obstruction vanishes arithmetically}
\label{sec:kt-global-mixed-cup}

The source reductions above do not make the remaining rational
negative root disappear.  Here the first mixed obstruction is
computed over the actual rational Galois group, with its ordinary
local correction.  The class-number unit removes this obstruction;
it does not force the negative class to be zero.  The construction
stops at a characteristic-$p$ Artin algebra of nilpotence degree
three and supplies no characteristic-zero source point.

\begin{lemma}[KT-AUXILIARY-CLASS-UNIT-MIXED-CUP]
Let $p>3$, let $D<0$ be a fundamental discriminant prime to $p$,
put $E=\mathbf Q(\sqrt D)$ and $\eta=\chi_D$, and set
\[
                  S=\{\infty,p\}\cup\{q:q\mid D\}.
                                                               \tag{MC.1}
\]
Let $k$ be a finite extension of $\mathbf F_p$.
If $B_{1,\eta}$ is a $p$-adic unit, then
\[
 H^2(G_{\mathbf Q,S},k(\eta))=0,\qquad
 \dim_k H^1(G_{\mathbf Q,S},k(\eta))=1,                  \tag{MC.2}
\]
and localization induces an isomorphism
\[
 H^1(G_{\mathbf Q,S},k(\eta))
       \ \xrightarrow{\ \sim\ }
 H^1(\mathbf Q_p,k(\eta))/H^1_{\rm unr}(\mathbf Q_p,k(\eta)).
                                                               \tag{MC.3}
\]
The finite set in (MC.1) is part of the assertion; no claim is
made after allowing arbitrary extra ramified primes.
\end{lemma}

\begin{proof}
It suffices first to take $k=\mathbf F_p$, and then extend
scalars: all cochain spaces and local maps are finite direct
sums after a finite residue-field extension.
Write $L=E(\mu_p)$ and $\Delta=\operatorname{Gal}(L/\mathbf Q)$.
The $a=0$ case of KT-AUXILIARY-UNIT-SELMER-BRIDGE gives
\[
              p\nmid h(E),\qquad
              \operatorname{Cl}(L)\{p\}^{\eta\omega}=0.
                                                               \tag{MC.4}
\]
For $M=\mathbf F_p(\eta)$ the Cartier dual is
$M^D=\mathbf F_p(\eta\omega)$.
A class in $\Sha^1_S(\mathbf Q,M^D)$ is unramified outside
$S$ and zero on each decomposition group in $S$.
Its restriction to $G_L$ is therefore an everywhere-unramified
homomorphism.  Since $|\Delta|=2(p-1)$ is prime to $p$,
restriction is injective and class field theory embeds this
group into
\[
   \operatorname{Hom}_\Delta
       (\operatorname{Cl}(L)/p,\mathbf F_p(\eta\omega)).  \tag{MC.5}
\]
The character in (MC.5) is the displayed character, not its
inverse: equivariance is $f(\sigma x)=\eta(\sigma)
\omega(\sigma)f(x)$.  Thus (MC.4) makes (MC.5) zero.

Poitou--Tate duality pairs this zero group perfectly with
$\Sha^2_S(\mathbf Q,M)$.  At $p$, local duality gives
$H^2(\mathbf Q_p,M)=H^0(\mathbf Q_p,M^D)^\vee=0$,
because $\omega$ is nontrivial on inertia and $\eta$ is
unramified.  At $q\mid D$, the character $\eta\omega$
is nontrivial on inertia because $q\ne p$.  The local
$H^2$ is again zero.  At the real place it is zero by
averaging, since $p$ is odd.  Hence the global $H^2$ equals
$\Sha^2_S$ and is zero.  Global Euler characteristic for
the odd character $\eta$ now gives
$h^0-h^1+h^2=-1$.  Since $h^0=h^2=0$, this proves (MC.2).

Local Euler characteristic at $p$ gives $h^1=h^0+1$.
If $\eta(p)=-1$, the two dimensions $h^1$ and
$h^1_{\rm unr}$ are $1,0$; if $\eta(p)=1$, they are $2,1$.
Thus the target of (MC.3) has dimension one in either case.
Its kernel is zero.  Indeed a class in that kernel restricts
to an everywhere-unramified homomorphism on $G_E$: it is
unramified at $p$ by the kernel condition; outside $S$ this
is built into $G_{\mathbf Q,S}$; and at every $q\mid D$,
local duality and local Euler characteristic give
\[
               H^1(\mathbf Q_q,\mathbf F_p(\eta))=0
                                                               \tag{MC.6}
\]
because both local $H^0$ and $H^2$ vanish.
The homomorphism factors through $\operatorname{Cl}(E)/p=0$
by (MC.4).  Restriction to $G_E$ is injective because its
index is two.  This proves injectivity, and equality of
dimensions proves (MC.3).
\end{proof}

\begin{theorem}[KT-GLOBAL-QUADRATIC-MIXED-PERSISTENCE]
Retain the preceding hypotheses, let $e\ge2$ be even with
$(p-1)\nmid2e$, and put
\[
        \psi=\omega^e,\qquad
        \alpha=\eta\psi,\qquad \beta=\psi^{-1}.
                                                               \tag{MC.7}
\]
Let $u$ be a continuous cocycle in
$Z^1(G_{\mathbf Q,S},k(\alpha))$.
Let $v$ be a continuous cocycle in $Z^1(G_{\mathbf Q,S},k(\beta))$
whose cohomology class is zero at $\mathbf Q_p$ and unramified
on inertia at every finite prime other than $p$.
There is a continuous polarized residual deformation over
\[
                           B=k[s,t]/(s^2,t^2)             \tag{MC.8}
\]
of
\[
 r_0(g)=(\operatorname{diag}(\psi(g),1,\psi(g)^{-1}),\eta(g))
\]
in $\mathrm{GL}_3(B)\rtimes\{\pm1\}$, with first-order
terms
\[
       s\,u(g)(E_{12}+E_{23})+
                         t\,v(g)(E_{21}-E_{32}).          \tag{MC.9}
\]
Its restriction to $G_E$ is unramified away from $p$.
At $p$, it is conjugate by a matrix congruent to one to an
upper-triangular family with the prescribed specialized
inertial diagonal $(\psi,1,\psi^{-1})$.
Thus the two directions are compatible through their mixed
quadratic term, with the actual ordinary local condition.
\end{theorem}

\begin{proof}
Set
\[
 X=E_{12}+E_{23},\qquad Y=E_{21}-E_{32},\qquad
 Z=[X,Y]=\operatorname{diag}(1,-2,1).                     \tag{MC.10}
\]
The unitary outer automorphism is
$\theta(g)=J(g^t)^{-1}J$, with $J$ the antidiagonal identity.
The adjoint action of $r_0$ on $X,Y,Z$ has characters
$\alpha,\beta,\eta$, respectively.
As $(s,t)^3=0$, logarithms and exponentials terminate after
degree two; division by two is permitted.
For an unknown continuous cochain $w:G_{\mathbf Q,S}\to k$,
put
\[
 L(g)=s u(g)X+t v(g)Y+st w(g)Z,
 \qquad r(g)=\exp(L(g))r_0(g).                            \tag{MC.11}
\]
The Baker--Campbell--Hausdorff identity here is a finite
polynomial identity, not a convergence assertion.
It says that $r$ is a homomorphism exactly when
\[
 \begin{split}
 w(gh)-w(g)-\eta(g)w(h)
   =\tfrac12\bigl(&u(g)\beta(g)v(h)\\
                 &-v(g)\alpha(g)u(h)\bigr).
 \end{split}                                             \tag{MC.12}
\]
The right side is one half of the difference of the two
cup-product cochains.  It is a continuous $2$-cocycle in
$k(\eta)$ because $\alpha\beta=\eta$.
By (MC.2) a continuous $w$ solving (MC.12) exists.
This constructs the global semidirect-product homomorphism.
All three matrices in (MC.10) have trace zero, so its left
matrix has determinant one.  The right component is exactly
$\eta$, preserving the stated polarization.

\paragraph{The ordinary correction at $p$.}
Since $[v]|_{G_{\mathbf Q_p}}=0$, choose $d\in k$ with
$v(g)=\beta(g)d-d$ on this decomposition group.
Conjugate (MC.11) by $\exp(tdY)$.
The negative first-order term becomes zero, and the logarithm
of the remaining deformation factor is
\[
       s u(g)X+st w_p(g)Z,\qquad
       w_p(g)=w(g)-\tfrac12d(1+\beta(g))u(g).             \tag{MC.13}
\]
There are no other terms because $s^2=t^2=0$.
The homomorphism property shows that $w_p$ is now a
$k(\eta)$-valued cocycle on $G_{\mathbf Q_p}$.
By (MC.3), a global cocycle $h$ in $k(\eta)$ can be chosen
so that $[w_p+h]$ is unramified at $p$.
Replacing $w$ by $w+h$ preserves (MC.12).
Because $\eta$ is trivial on inertia at $p$, coboundaries
vanish on that inertia.  Unramifiedness of the class therefore
means that the actual cochain $w_p+h$ is zero on inertia.
Formula (MC.13) is upper triangular and its diagonal
correction on inertia is now zero.  This is precisely the
claimed ordered ordinary family.  Restriction to $G_E$
and conjugation transport the flag to either place above $p$.
The argument also works when $p$ is inert.

\paragraph{Away-prime inertia.}
There is no ramification outside $S$ by construction.
At $q\mid D$, both $\alpha$ and $\eta$ are nontrivial on
inertia, while the cyclotomic character is unramified.
Local duality and local Euler characteristic give
\[
 H^1(\mathbf Q_q,k(\alpha))=0,
 \qquad H^1(\mathbf Q_q,k(\eta))=0.                     \tag{MC.14}
\]
First conjugate by a matrix $\exp(saX)$ to remove the local
$X$-cocycle.  The remaining logarithm has the form
$t v(g)Y+st w_q(g)Z$, with $w_q$ an $\eta$-cocycle.
The second vanishing in (MC.14) lets a further conjugation
by $\exp(stbZ)$ remove that cocycle.
The coefficient of $Y$ is unramified by hypothesis on $v$.
On $I_q\cap G_E$ the resulting representation is therefore
the identity, as is $r_0$.  Conjugating back preserves this
identity.  This proves the asserted away-prime property.
\end{proof}

\paragraph{What the calculation decides.}
The odd character $\alpha=\eta\omega^e$ has
$h^1-h^2=1$ by global Euler characteristic, so a nonzero
positive class $[u]$ always exists at the set (MC.1).
A nonzero strict negative class $[v]$, in contrast, is an
instance of the original rational strict Selmer problem.
The theorem says that, \emph{if that class exists}, the
class-number unit supplies no contradiction from the first
mixed cup product or its ordinary local correction.
It does not assume that class zero, and it does not prove
that it exists.

The family (MC.8) has characteristic $p$ and nilpotents.
No lift through all higher orders, no characteristic-zero
irreducibility, no automorphy, and no point of the normalized
global source ring has been constructed.  At the local
ordinary functor it is a genuine Borel family followed by
near-identity conjugation, as in the already checked
KT-ORDINARY-INTEGRAL-COMPONENT.  Its noncritical residual-character
condition holds here: every nonzero difference among $e,0,-e$
is even, so none is congruent to $1$ modulo the even integer
$p-1$.  No characteristic-zero crystalline-point condition
on the Artin ring $B$ is being imposed.
This verifies the local condition but does not resolve the
higher global equations.
Accordingly this section is an exact obstruction audit,
not an additional arithmetic vanishing gate.

\paragraph{Imported duality and exact source scope.}
We use J.~S.~Milne, \emph{Arithmetic Duality Theorems},
second edition (2006), author PDF
\url{https://www.jmilne.org/math/Books/ADTnot.pdf}.
The definitions of $G_{\mathbf Q,S}$ and $\Sha^r_S$ are in
Chapter~I, pp.~48 and 56; the finite-module Poitou--Tate
pairing is Theorem~4.10(a), pp.~56--57.
Local duality is Corollary~2.3, p.~28, and local Euler
characteristic is Theorem~2.8, p.~31.
Global Euler characteristic is Theorem~5.1, p.~67, with
the finite-$S$ setup on p.~66 and its useful cohomological
form restated in Remark~5.2(c), p.~68.
These complete statements and surrounding definitions were
read for the present application.  Their foundational proofs
are imported as theorems; no new proof of Poitou--Tate duality
or arithmetic comparison with etale cohomology is claimed.
